Erklären sie die wesentlichen Kerngrößen der ersten Säule von Solvency II


Zu den wesentlichen Kerngrößen der ersten Säule zählt der Risikokapitalbedarf, dieser wird differenziert in 
- Solvabilitätskapitalanforderung (SCR)
- Mindestkapitalanforderung (MCR)

Das SCR enspricht dem VaR der Basiseigenmittel über den Zeitraum eines Jahres des VU ,
zu einem Konfidenzniveau von 99,5\%.
Das bedeutet, dass die Verluste der Eigenmittel
diesen VaR-Wert mit einer Wahrscheinlichkeit von 99,5\% innerhalb eines Jahres nicht übersteigen.

Das MCR stellt die Untergrenze der anrechnungsfähigen
Eigenmittel dar, bei deren Unterschreitung die Aufsichtsbehörde einen kurzfristigen und
realistischen Finanzierungsplan einfordert (oder sogar die Geschäftserlaubnis entziehen kann).

Die Solvabilitätskapitalanforderung kann zum einen mit der Standardformel bewertet
werden, bei der nur bestimmte Parameter unternehmensspezifisch gesetzt werden können.
Zum anderen kann ein (partielles) internes Modell gewählt werden.

Die Bewertung der Bilanz verläuft anders als nach HGB und wird auch Solvenzbilanz gennant:

Den größten Posten auf der Passivseite bilden die versicherungstechnischen Rückstellungen. 
Diese bestehen aus der Summe des besten Schätzwerts der versicherungstechnischen Verpflichtungen und einer
Risikomarge, die jeweils unter Einhaltung der zugrundeliegenden Vorschriften gesondert
ermittelt werden.

Die Bewertung der Kapitalanlagen auf der Aktivseite der Solvabilitäts-
übersicht erfolgt zu Marktwerten.

INSERT IMAGE PAGE 86


Erläutern sie die ökonomische Bewertungsmethoden in der Solvenzbilanz

für die Bewertung die Positionen der Solvenzbilanz unter dem Solvency
II-Regime folgende Bewertungsansätze:
• Aktiva: Fair Value/Mark to Market
• Passiva: Fair Value/Mark to Model
• Versicherungstechnische Rückstellungen:
+ Best Estimate
+ Risk Margin

Grundsätzlich gibt es zwei unterschiedliche Arten von Verbindlichkeiten:
1. Die absicherbaren (hedgebaren) und
2. die nicht absicherbaren (nicht hedgebaren) Verbindlichkeiten

Absicherbare Verbindlichkeiten sind jene, die im Zusammenhang mit Versicherungsverträgen zu
berechenbaren Zahlungsströmen führen. Es lässt sich die Mark to Market-Methode anwenden.

Für nicht absicherbare (nicht hedgebare) Verbindlichkeiten, wie die versicherungstechnischen
Rückstellungen, liegt kein Marktwert vor.
Sie werden mithilfe mathematischer Modelle (Mark to Model) geschätzt (bester Schätzwert = Best Estimate, Mark to Model).
Für die Unsicherheit des besten Schätzwertes ist zusätzlich noch eine Risikomarge zu ermitteln.

Im Allgemeinen werden Cashflow-Projektionen eingesetzt, also gibt es keine Abgrenzungen/Beitragsüberträge.


Erklären sie anhand des Beispiels von Schaden/Unfallversicherungen die Best Estimate Methode

\(\text{Best Estimate} = \text{Schadenrückstellung} + \text{Prämienrückstellungen} \)

Schadenrückstellungen(IBNR) sind für bereits entstandene aber noch nicht eindeutiv evaluierte Schäden,
es wird differenziert in:
BNYR = Incurred but not yet reported
IBNER = Incurred but not enough reported
IBNR = Incurred but not reported = IBNYR + IBNER

für diese werden Techniken wie Chain Ladder, additive Methoden genutzt.

Prämienrückstellungen sind für etwaige Schäden aus bereits
eingegangenen vertraglichen Verpflichtungen zu bilden.
Für jedes einzelne Risiko (Versicherungsvertrag) ist sowohl das
Auftreten, als auch die Höhe des auftretenden Schadens zu bewerten.

Für die Prämienrückstellungen wird als Näherungsformel von EIOPA empfohlen

\( \text{BE} = \text{CR}\cdot\text{VM} + (\text{CR} -1)\cdot \text{PVFP} + \text{ARE}\cdot \text{PVFP} \)
BE := bester Schätzwert der Prämienrückstellung.
CR := Combined Ratio d.h. (Schäden + schadenbezogene Kosten) / (verdiente Prämien)
VM := Beitragsübertrag: Prämien für laufende Verträge, abzüglich der Prämien, die bei diesen Verträgen
PVFP := Barwert künftiger Prämien.
ARE := Schätzung der Abschlusskostenquote.


Was ist die Risikomarge (risk margin) und wie wird Sie berechnet?

Für hedgebare Risiken ist keine Risikomarge vorzuhalten.

Für nicht hedgebare Risiken ist eine Risikomarge zu ermitteln. Diese Risikomarge ergibt sich aus
der Unsicherheit des besten Schätzwertes (Best Estimate) für die nicht hedgebar Risiken.
Sie ist der Betrag, den ein anderes (risikoaverses) Versicherungsunternehmen (Referenz) über den Best Estimate 
der versicherungstechnischen Rückstellungen hinaus fordern
würde, um die Versicherungsverpflichtungen zu übernehmen (exit value). Für die Berechnung
der Risikomarge wird das Kapitalkostenverfahren (Cost of Captial = CoC) angewendet

Die Schritte zur Rechnung sind:
1. Projektion der Solvenzkapitalanforderung (SCR) für die Zukunft bis zum Auslaufen
   der jetzigen Verpflichtung in \( t = 0 \) je Sparte \( \text{SCR}_i \):
2. Ermittlung der Kapitalkosten (CoC), die für das Vorhalten der im 1. Schritt ermittelten
   zukünftigen Solvenzkapitalanforderungen notwendig sind. Darauf ist der CoC/Kapitalkostensatz als Risikomarge
   einzurechnen (typischerweise 6\%).
3. Diskontierung der Zahlungsreihe

\( \text{RM} = \text{COC}\cdot \sum_{i=1}^m\sum_{t=0}^n \frac{1}{(1 + r_{t+1})^{t+1}}\cdot \text{SCR}_i(t) \)

es ist 
\(r_t\) := risikoneutraler Zinssatz im Abwicklungsjahr t
\( n \) := letztes Abwicklungsjahr
\( m \) := Anzahl der Sparten

Was ist Kapitalallokation und welche Rolle nimmt sie bezüglich der SCR Kalkulation ein?

Werden unterschiedliche Risiken auf übergeordneter Ebene zusammengefasst, so addieren sich die benötigten Risikokapitalien nicht, 
sondern müssen entsprechend der Abhängigkeitsstruktur aggregiert werden.

In der Regel entsteht ein Diversifikationseffekt, d.h. das aggregierte benötigte Risikokapital auf übergeordneter Ebene ist geringer 
als die Summe der einzelnen Kapitalien. Dies spiegelt sich in der subadditivität von Risikomaßen wieder
\( \rho(\sum_i X_i) \leq \sum_i \rho(X_i) \)


Definition: Kapitalallokation 
Seien \( X_{1}, ..., X_{m} \) Zufallsgrößen, die \( m \) Portfolios beschreiben und \( X = \Sigma_{i =1}^{m} X_{i} \) das Gesamtrisiko.
Eine Kapitalallokation \( K = (K_{1}, ..., K_{m} \) heißt Zuteilung für das Risikomaß \( \rho \), wenn das gesamte Kapital allokiert wird,d.h. 
\( \rho(X) = K_{1} + K_{m} \) und das allokierte Kapital nicht höher als das entsprechende Einzelkapital ist, d.h. 
\( K_{i} \leqslant \rho(X_{i}), i = 1, ..., m \).


Geben Sie Beispiele von Kapitalallokationsprinzipien


Das Verfahren der proportionalen Aufteilung ist einfach, hat jedoch den Nachteil, 
die Korrelationen der Einzelportfolios nicht zu berücksichtigen.

\( K_{i} = \frac{\rho(X_{i})}{\Sigma^{m}_{i = 1}\rho(x_{i})} \cdot \rho(X) \).


Definition: Inkrementelle Allokation
Nach dem diskreten Marginalprinzip von Merton / Perolderfolgt(???) die Aufteilung proportional zum Beitrag des Geschäftsbereichs zum Gesamtrisiko. 
Ziel ist es für die bestehende Aufteilung in Geschäftsbereiche das Kapital möglichst gerecht zuzuordnen.

\( K_{i} = \frac{\rho(X) - \rho(X - X_{i})}{\Sigma^{m}_{i = 1}(\rho(X) - \rho(X - X_{i}))} \cdot \rho(X) \)


Definition: Kovarianzbasierte Allokation 
Im Spezialfall, dass das Risikomaß  \( \rho = \beta\sigma \) ein Vielfaches der Standardabweichung  \( \sigma \) ist, 
erhalten wir das Kovarianzprinzip.

\( K_{i} = \beta \cdot \frac{Cov(X_{i}, X)}{\sigma} \)

Rechnen Sie ein konkretes Beispiel zur Kapitalallokation vor mit einer Methode ihrer Wahl

VU mit drei Sparten (Hausrat \( X_{1} \), Haftpflicht \( X_{2} \) und Wohngebäude \( X_{3} \)).
\( X_i \sim \mathcal{N}(0, \sigma_i^2) \), \( \sigma_i \), \( \sigma^{2}_{1} = 100 \), \( \sigma^{2}_{2} = 49 \),  \( \sigma^{2}_{3} = 25 \)
und Korrelationsmatrix  \( \left(\begin{array}{ccc}1 & 0 & 0,6 \\0 & 1 & 0 \\0,6 & 0 & 1\end{array}\right) \)

\( \sigma_{X} = \sqrt{\sigma^{2}_{1} + \sigma^{2}_{2} + \sigma^{2}_{3} + 2\cdot0,6\cdot\sigma_{1}\cdot\sigma_{3}} = \sqrt{100 + 49 + 25 + 2\cdot0,6\cdot\sqrt{100}\cdot\sqrt{25}} = \sqrt{234} = 15,3 \)


Wir nutzen das proportionale Prinzip mit \( \rho = 1,25\cdot\sigma(x) \)

\( \rho(X) = 1,25 \cdot 15,3 = 19,1 \quad \rho(X_{1}) = 1.25\cdot\sqrt{100} = 12,5 \quad \rho(X_{2}) = 8,75 \quad \rho(X_{3}) = 6,25 \)
\( K_{i} = \frac{\rho(X_{i})}{\Sigma^{m}_{i = 1}\rho(x_{i})} \cdot \rho(X) \).


\( K_{1} = \frac{\rho(X_{1}}{\rho(X_{1}) + \rho(X_{2}) + \rho(X_{3})} \cdot \rho(X) = \frac{12,5}{12,5 + 8,75 + 6,25} \cdot 19,1 = 8,69 \)
\(K_{2} = \frac{8,75}{27,5} \cdot 19,1 = 6,085 \)
\(K_{3} = \frac{6,25}{27,5} \cdot 19,1 = 4,35 \)




